\section{Existence: constructions, the Schur/Perov machinery, and the long-run law}
\label{app:existencemachinery}

This appendix collects the technical constructions deferred from
\Cref{sec:existence,sec:assumptions}: the alternative existence
routes, the Schur-condition matrix computations, the purified
(shock) constructions and the invariant-law machinery,
and the discharge apparatus for the assumption residual clauses.
None of it changes any statement in the main text.

\paragraph{The coupling matrix.} With $L_{TS}$, $L_{ST}$ and
$\beta_{S,\mathrm{row}}$ as in \Cref{thm:existence} below,
\[
\mathbf M=
\begin{pmatrix}
\beta_T & L_{TS}\\
L_{ST} & \beta_{S,\mathrm{row}}
\end{pmatrix},
\]
and $\rho(\mathbf M)$ is its spectral radius. When the four coefficients are
non-negative and the two diagonal ones are below one, the boxed Schur
inequality of \Cref{thm:existence} is equivalent to $\rho(\mathbf M)<1$
(\path{specRadius_lt_one_iff_schurDet}); the theorem itself imposes no sign on
the coefficients.

The theorem, which \Cref{sec:existence:general} states in words, is the
following.

\begin{theorem}[Existence on invariant sets under the Schur condition]
\label{thm:existence}
Fix closed, non-empty sets of value functions $\mathcal C_T$ and
$\mathcal C_S$. Suppose that, for every sender value in $\mathcal C_S$, the
sender's dynamic argmax is unique at every state and every feasible compliance
level $a\in[0,1]$, and that the target argmax is unique at every value pair from the two sets.
Together with the maintained continuity of the primitives, these
conditions define continuous feasible selectors. Suppose also that the
joint Bellman operator maps
$\mathcal C_T\times\mathcal C_S$ into itself. Denote its two rows by
$\mathcal T_T$ and $\mathcal T_S$. Suppose, for all value pairs
$z=(\tilde V_T,\tilde V_S)$ and $z'=(\tilde V_T',\tilde V_S')$ from the two sets, that
\begin{align}
\|\mathcal T_T(z)-\mathcal T_T(z')\|_\infty
&\leq
\beta_T\|\tilde V_T-\tilde V_T'\|_\infty+L_{TS}\|\tilde V_S-\tilde V_S'\|_\infty,\\
\|\mathcal T_S(z)-\mathcal T_S(z')\|_\infty
&\leq
L_{ST}\|\tilde V_T-\tilde V_T'\|_\infty+
\beta_{S,\mathrm{row}}\|\tilde V_S-\tilde V_S'\|_\infty,
\end{align}
where $\beta_T,\beta_{S,\mathrm{row}}<1$ and
\begin{equation}
\boxed{
L_{TS}L_{ST}
<
(1-\beta_T)(1-\beta_{S,\mathrm{row}})
}
\label{eq:Schur}
\end{equation}
(no sign is imposed on the four coefficients; when they are non-negative,
\eqref{eq:Schur} is equivalent to $\rho(\mathbf M)<1$). Then the joint
Bellman operator has a unique fixed point on the invariant product set in a weighted
sup-norm. The associated optimal target policy and greatest sender response
form an equilibrium $E=(\Vt,\Vs,\alpha)\in\mathcal E$ with
$(\Vt,\Vs)\in\mathcal C_T\times\mathcal C_S$ and $\alpha=\alpha_{\Vt,\Vs}$,
the target selector the two sets define. Every equilibrium
$E'=(V_T',V_S',\alpha')\in\mathcal E$ with
$(V_T',V_S')\in\mathcal C_T\times\mathcal C_S$ and
$\alpha'=\alpha_{V_T',V_S'}$ coincides with $E$: $V_T'=\Vt$, $V_S'=\Vs$ and
$\alpha'=\alpha$.
\end{theorem}

A negative coefficient is allowed because it can only describe a degenerate
invariant set. If $\beta_T<0$ or $L_{ST}<0$, the row bound at a fixed sender value
forces the target set to be a single point, and symmetrically
$L_{TS}<0$ or $\beta_{S,\mathrm{row}}<0$ makes the sender set a single
point; the bounds then hold again with non-negative coefficients, and the
Perov step applies unchanged. Both diagonal bounds are assumed in exchange:
with non-negative coefficients $\beta_T<1$ would follow from the other
conditions, but without the signs it does not. The coincidence clause is
uniqueness within the stated invariant sets and selector class; it is not global
equilibrium uniqueness (\path{feedbackStackelbergTop_unique_in_carrier_schur}).
The sender's argmax need be unique only at compliance levels in $[0,1]$: the
target chooses compliance there, and the target's Bellman operator reads the
sender's response only there (\path{Params.DynamicArgmaxUniqueOnI}). Uniqueness
at every real compliance level implies this (\path{DynamicArgmaxUnique.toOnI}),
and the corner economies of \Cref{prop:primitive-persistence} obtain the weaker
form with no bound on the curvature of their transitions
(\path{SymbolicEconomyF.DAUOnI_on_cone}).

\subsection{Alternative existence routes}

\Cref{thm:existence} uses a model-specific Schur-condition
row-Lipschitz route, with the invariant set of state-Lipschitz value functions
motivated by \citet{Curtat1996}. The literature offers several parallel sufficiency routes: Tarski
monotone iteration on a lattice of value pairs
\citep{BalbusReffettWozny2014JET,Vives1990}, Schauder on a sup-norm-compact
Lipschitz ball \citep{RaadWozny2019}, measurable selection on a decomposable
kernel \citep{HeSun2017}, and the fixed-point route of the risk-sensitive
overlapping-generations model \citep{JaskiewiczNowak2014}, which is a Glicksberg argument rather than a row-Lipschitz one. None is adopted as primary
here: the lattice route needs an order-preservation structure the coupled
operator does not have (\Cref{rem:twisted-order}); the Schauder and
measurable-selection routes return existence without the greatest-selection
convention or the closed-form constants the comparative statics use; and the
risk-sensitive extension is not needed for time-additive preferences. The
Schur-condition route is described in detail below. The general taxonomy of stationary
Stackelberg-equilibrium concepts for discounted stochastic games is
in \citet{Bucarey2023}, who formalize the stationary strong
Stackelberg equilibrium (SSSE) and fixed-point equilibrium (FPE)
notions, prove their existence and coincidence for classes of games
including myopic follower strategy and team games, and provide
counter\-examples where existence fails in the general case. The
restricted assumption set
$\{$\Cref{ass:sender},\Cref{ass:target},\Cref{ass:concavity}$\}$
describes structural inputs; the Schur-condition construction of
\Cref{thm:existence} also requires the value-set, selector, invariance
and row-bound hypotheses stated there, and the construction of
\Cref{prop:placable-family} discharges its own conditions from its fields.
Neither makes the three assumptions alone a sufficient condition for
existence.

\begin{remark}[Why the mixed-monotone order fails, and where]
\label{rem:twisted-order}
The lattice route is usually attempted twice: first under the pointwise product
order on value pairs, and then, when that fails, under the \emph{twisted}
order that reverses one coordinate,
\[
  (V_T,V_S)\preceq_{\mathrm{tw}}(\tilde V_T,\tilde V_S)
  \quad\Longleftrightarrow\quad
  V_T\le \tilde V_T \ \text{ and }\ \tilde V_S\le V_S \ \text{ pointwise.}
\]
The second attempt is the substantive one, because our pointwise refutations
vary the target continuation with the sender continuation held fixed and so
leave the twisted order untouched. Isotonicity of the joint Bellman map under
$\preceq_{\mathrm{tw}}$ requires four coordinate signs. The target row's
own-value sign holds, from $\beta_T\ge0$ alone. Its cross sign does not.

The mechanism of the failure is worth stating, because it is not an artifact.
The sender's continuation value enters the target's Bellman integrand only
through the sender's best response. Raising that continuation makes the sender
sanction more heavily, which costs the target current payoff but \emph{raises}
its continuation value, because a heavier sanction raises next period's
entrenchment and the target's value is increasing in entrenchment. Writing $L_T$
for the target's stage loss per unit of additional sanction, $g_T$ for the lower
slope of its value in the state, and $d_S$ for the lower bound on the transition's
sanction slope, the second effect dominates whenever
\[
  L_T \;<\; \beta_T\,g_T\,d_S ,
\]
which is the strict form of the entrenchment condition
$c_T\le\beta_T\,g_T\,d_{\mathrm{lo}}$ maintained throughout
\Cref{sec:existence}. That condition is the one under which the target
tolerates sanctions rather than complying in the persistence certificate.

The lattice route therefore fails on the economies this paper is about. The
refutation is a configuration: two ordered sender continuation values between
which the sender's optimal response strictly shifts, a target maximizer, the
difference-quotient inequalities and the strict entrenchment condition above
(\path{not_antitone_T_T_joint_in_sender_value}). It defeats antitonicity of
the target row on any domain containing that configuration, which is what the
unrestricted twisted-order route needs. Three scope points. It does not show
that every smaller value set fails: a value set on which the entrenchment condition is reversed, or
on which the sender's response does not move between the two values (a constant
shift of the continuation adds an action-independent term and moves nothing),
is not addressed. It closes a proof \emph{route}; existence is delivered by
the contraction argument of \Cref{thm:existence} and is untouched. And the
refutation is empty when
$\beta_S=0$, where the sender is myopic, its response does not depend on its own
continuation, and the target operator is constant in $V_S$, so the obstruction
is a genuine consequence of two-sided forward-looking behavior rather than of
the sequential timing.
\end{remark}


\subsection{The construction behind the theorem, in brief}
\label{app:jointiteration}

The proof of \Cref{thm:existence} is a joint Banach iteration on the value
pair in an invariant set with an asymmetric weighted sup-norm, whose contraction modulus is
governed by the Schur condition \eqref{eq:Schur}: the set is complete, two
invariance conditions keep the operator inside it, the row estimates have the
entries of $\mathbf M$ as coefficients, and the Perov $2\times2$ fixed point
(\citealp{Toda2021}; the weighted-max-norm reduction of
\citealp[\S2]{BertsekasTsitsiklis1996}) yields the value pair from which the
optimal target policy and the greatest sender response are read off. Three
levels of the same route are kept apart in the formalization (the
conditions at a supplied equilibrium, at the constructed fixed point, and
over the whole invariant set), and the strongest is the checkable \emph{ex ante}
certificate (\path{Sanctions/Theory/Equilibrium/PerovExistence.lean},
\path{Sanctions/Theory/Certificates/TwoSidedRuleComparison.lean};
\Cref{app:formalization}). The fixed point is locally isolated inside its
invariant set: a shape-conforming equilibrium within an enlarged norm ball at
which the Schur condition still holds coincides with it. This is a local
uniqueness statement inside a specified invariant set, with no genericity,
manifold, nondegeneracy or finiteness claim about the global equilibrium
set; the smallest invariant radii of the balls are
$\|\Ut\|_\infty/(1-\beta_T)$ and $\|\Us\|_\infty/(1-\beta_S)$. The anatomy
of the iteration, the isolation statement and the coupled low-$\delta$
diagnostics are not reproduced in this appendix.

\subsection{Contraction restoration under the shock}

\paragraph{Machine-checked contraction restoration under the shock.}
The purified existence above is, in the cited references, a
\emph{topological} fixed point (a continuous self-map on a compact value
class, closed by Fan--Glicksberg). A sharper statement of the mechanism on the sender side follows (the construction is mapped in \Cref{fig:shock-construction}). Model
the atomless shock as a smoothing operator on continuation values,
$(\mathrm{smear}\,V)(y)=\mathbb{E}_\eta[V(y+\eta)]$ with $\eta$ drawn from the shock
density. Two properties of $\mathrm{smear}$ are exactly what an atomless
smooth density delivers: it is sup-norm \emph{nonexpansive} (averaging
against a probability density, $\int\kappa=1$), and it maps any bounded $V$
to a state-\emph{Lipschitz} function (mollifier smoothing); only the first is
carried by the Lean structure \path{ShockedParams}, whose single proof field is
nonexpansiveness. The first
property alone shows, with no new regularity, that the shocked sender
Bellman operator $V\mapsto\mathrm{bellmanS}(\alpha,\mathrm{smear}\,V)$
\emph{remains a $\beta_S$-contraction}: the deterministic
H\"older-$\tfrac12$ obstruction is bypassed because the shock enters only
through the nonexpansive continuation, so the contraction modulus is
unchanged at $\beta_S<1$, and the sender value is \emph{constructed} by the
Banach engine rather than postulated.
This is a Banach, rather than Fan--Glicksberg, existence route under
the shock. Nor is $\mathrm{smear}$ left abstract: the He--Sun shock is
carried as a concrete \emph{atomless probability kernel} $\kappa$ on the
transition (the He--Sun idiosyncratic component), with
$(\mathrm{smear}\,V)(y)=\int V(\mathrm{clip}(y+\eta))\,\kappa(\mathrm d\eta)$,
and its nonexpansiveness is \emph{derived} from $\int\kappa=1$ rather than
assumed, so the shocked corner equilibrium holds from a genuine kernel
primitive. The shocked equilibrium is an equilibrium of the \emph{shocked} game (its
Bellman equations carry $\mathrm{smear}$), a sibling object to the
deterministic \Cref{thm:corner-existence}, and it too now holds: the
shocked equilibrium structure, both shocked Bellman fixed-point equations
(sender and target), and a non-vacuous shocked \emph{corner} instance are
established, with the two value functions \emph{constructed} as Banach
fixed points of the shocked sender and target operators. The \emph{interior} (non-corner, state-varying leader) case is also now
established: the slaved shocked operator $\Phi(\Vs)=\mathrm{bellmanS\_shock}
(\mathrm{lead}(\Vs),\Vs)$ has modulus $\le\beta_S+L_{U_S,a}\,L_\alpha$, where the
continuation term contributes $\beta_S$ (nonexpansive $\mathrm{smear}$) and
the policy term contributes $L_{U_S,a}\,L_\alpha$ ($U_S$ Lipschitz in the
target action, leader policy $L_\alpha$-Lipschitz in $\Vs$); under the condition
$\beta_S+L_{U_S,a}\,L_\alpha<1$ this is a Banach contraction, and its fixed point
yields an interior shocked equilibrium with a state-varying leader policy. The
contraction-specific quantity is $L_\alpha$: deterministically the leader
argmax is only H\"older-$\tfrac12$ in $\Vs$, but the atomless shock's
mollifier-Lipschitz smoothing upgrades it to Lipschitz (the
$\mathrm{smear}$-regularity property), which the numerics confirm in the
relevant range of shocks.  In the canonical theorem it is derived from composite
strong concavity and the Stackelberg-coupling bound; joint follower continuity
supplies the objective continuity and the canonical leader selector.
The full-width case $h=1$ now discharges that follower-continuity input: the moving
average is exactly affine on the physical state interval for every bounded
continuous continuation, which yields a unique sender argmax and a jointly
continuous greatest selection on $[0,1]\times\mathbb R$.
At the $P_3$ economy, the same affine identity permits a stronger exact
construction that bypasses the reusable leader-regularity theorem altogether.
The two smeared continuation values are summarized by their endpoint spreads;
eliminating the target spread gives a quartic equation for the sender spread.
Lean isolates the economically relevant root in an exact rational interval,
constructs the induced piecewise-affine leader policy and clipped greatest
sender response, builds both bounded-continuous value functions, and verifies
global target optimality and both Bellman fixed-point equations through every
clipping region. The resulting full-width shocked equilibrium has no
strong-concavity, cross-value/action coupling, invariant-set, or policy
hypothesis.
The \emph{mechanism} of that upgrade is itself now
established: a general lemma shows that for a $c$-strongly-concave objective
whose two-maximizer difference is action-Lipschitz with constant $L$, the
maximizers satisfy $c\,|a_f-a_g|\le 2L$: a Lipschitz (not Hölder-$\tfrac12$)
argmax bound. The
deterministic obstruction bounds the same difference by its \emph{sup}
($\to\sqrt{}$, Hölder-$\tfrac12$); the smear-smoothed continuation makes it
\emph{action-Lipschitz} with constant $\propto\|V-V'\|$ (the
$\mathrm{smear\_lipschitz}$ property), turning the same inequality into the
Lipschitz rate. That $\mathrm{smear\_lipschitz}$ estimate is itself now
discharged from a concrete atomless (uniform) kernel: the moving-average
smoothing $(\mathrm{smear}\,U)(y)=\tfrac1{2h}\!\int_{y-h}^{y+h}U(\mathrm{clip}\,z)\,dz$
is state-Lipschitz with constant $\|U\|/h$, by the elementary window-shift bound
$|\!\int_{y-h}^{y+h}g-\!\int_{y'-h}^{y'+h}g|\le 2\|g\|_\infty|y-y'|$. The
argmax-core reduces $\mathrm{lead\_lip}$ to a single action-direction estimate
on the leader objective \emph{difference}. The
shock's $\sigma$- and $\Vt^{\mathrm{dyn}}$-Lipschitz bounds control
the \emph{magnitude} of that estimate, but its variation \emph{in the leader's
own action} is not supplied by the shock: because the sender response
$\sigma(x,a)$ enters the leader objective and depends on $a$, this
action-Lipschitz estimate (the Stackelberg coupling, isolated as the
hypothesis $\mathrm{h\_obj\_diff}$) is \emph{consumed}, not derived. A
reference sweep confirms this is not an artifact of our particular kernel:
smoothing the continuation regularizes the \emph{state} direction only (a
uniform-Lipschitz-in-action integrand need not be jointly Lipschitz;
\citealp{Hinderer2005}). Concretely, the moving-average smear is only $C^1$
(its second derivative would require a $C^1$ continuation, which value functions
are not), so the action-direction cross-derivative an implicit-function / envelope
derivation \citep[the Lipschitz-policy route of][]{Montrucchio1987} would need
does not exist. The feedback-Stackelberg constructions we hold do not
derive the analogous leader-side regularity from primitives either; where
they need it, it is imposed. So
$\mathrm{lead\_lip}$ is the argmax-core applied to one consumed coupling
estimate. The canonical interior theorem pins the leader objective to the
genuine target payoff $\mathrm{joint\_T\_top\_integrand}$.  Joint continuity
of the shocked sender selector and strong concavity of this real composite
objective at any positive modulus \(c_T\) then produce a unique canonical
leader maximizer; Berge continuity, feasibility, and pointwise leader
optimality are proved rather than separately assumed.  For a general
half-width, the remaining analytic surface is therefore exactly those two
conditions plus the coupling estimate and the direct contraction condition
\(\beta_S+2L_{\mathrm{lead}}/c_T<1\), not an externally supplied leader map or
a fixed numerical curvature.  In the state-local full-width specialization,
follower continuity is proved and the remaining surface contracts to composite
leader strong concavity, the coupling estimate, and this same symbolic condition.
The kernel side is now also inhabited by a \emph{concrete}
uniform-kernel parameter, and the smear's state-Lipschitz constant is the closed
form $\|V\|/h$ rather than a named primitive. This does not discharge, or prove
joint satisfiability of, the general three-condition analytic package. The
conditional theorem is wired to the mechanism-active $P_3$ economy: once that
package is supplied, it builds a purified equilibrium as a Banach fixed point
with no Schur condition. Its abstract state-local $h=1$ sibling discharges follower
continuity but still consumes composite strong concavity, the coupling estimate,
and \(\beta_S+2L_{\mathrm{lead}}/c_T<1\). The older modulus-\(2\) and
\(L_{\mathrm{lead}}<1/2\) statement is only the arithmetic specialization at
the printed example parameter vector. These conditional templates are not headline existence
results.

The exact $h=1$ construction is logically separate and stronger at that
parameter vector. It solves the equilibrium rather than proving the abstract
leader hypotheses over an operator domain. The selected sender endpoint spread
is an exact root of a quartic in a rational bracket; the target spread is then
recovered algebraically. These spreads determine a feasible continuous
piecewise-affine leader policy, an exact clipped greatest sender response, and
both bounded-continuous values. Lean verifies the sender fixed point, the target
fixed point, and global target optimality through the lower-clipped, interior,
and upper-clipped sender regions. Hence an unconditional, genuinely
state-varying full-width $P_3$ equilibrium is exhibited.

The dynamic-over-myopic theorem still requires its separate basin inputs.
At full width, monotonicity of the sanction rule is available twice. The reusable
canonical-selector result follows from the change of variable
$b=a-(27/10)x$, strict decreasing differences of the target objective, and
interval dominance of the sender objective. At the explicit equilibrium there
is also a direct proof: the leader policy is increasing, the sender vertex
remains increasing after the leader response is substituted, and clipping
preserves monotonicity. Thus the equilibrium's own sanction rule, under the greatest selection, is
monotone on all of $[0,1]$.

The separate dynamic floor is nevertheless impossible. The full-width sender
selector is the clip of
$\tfrac98x-\tfrac1{12}-\tfrac5{12}\alpha(x)
 +(V_S(1)-V_S(0))/24$, and the top-state floor holds if and only if
$V_S(1)-V_S(0)\ge10\alpha(1)-1$. The sender Bellman equation bounds the same
spread above by
$\tfrac58(\alpha(1)-\alpha(0))+\tfrac{35}{16}$, so a putative floor forces
$\alpha(1)\le17/50$. The target Bellman equation gives
$V_T(1)-V_T(0)\le60/11$. Combining these restrictions with the exact sender
vertex makes a deviation at $x=1$ to $a'=1/2$ strictly profitable, contradicting
target optimality. Hence no full-width shocked equilibrium has the top-state
floor, and no non-empty basin ending at $1$ satisfies the full-width floor
component. In particular, the explicit equilibrium has a monotone sanction
rule and a top-state sanction strictly below one. It therefore cannot be
assembled with the floor into full pathwise ordering at $P_3$; the myopic
basin inputs also remain open. The earlier test continuation $V_S(x)=-100x$
remains a useful disproof-first check, while the equilibrium argument closes the
stronger fixed-point statement.

For a general half-width, the three-condition analytic package remains the
reusable theorem's honest assumption surface. The stage curvature
$\partial^2U_T/\partial a^2=-2$ does not by itself establish strong concavity
of the composite leader objective, because the target action moves the clipped
sender response and hence the continuation. An exact Lean counterexample at
$x=5/12$ shows that the static composite already violates $2$-strong concavity
on the rational actions $9/10$, $37/40$, and $19/20$. Likewise, smoothing
regularizes the state direction, not the leader-action direction controlled by
the Stackelberg-coupling estimate. These facts explain why the abstract theorem
retains those inputs; they are not hypotheses of the explicit full-width
equilibrium. The deterministic model is the $\kappa\to\delta_0$ limit.

\subsection{The derived dividing line, in general form}
\label{app:derived-dividing-line}

\Cref{prop:derived-dividing-line} is stated for the equilibrium closed loop on
the unit state space. The content is an order-topological fact about a
continuous self-map, and it is proved at that level; the equilibrium statement
substitutes $G=G_E$ and nothing else.

\begin{proposition}[Derived crossings of a continuous self-map]
\label{prop:derived-crossings}
Let $G$ be continuous on $[a,b]$ with $G(a)=a$ and $G(b)=b$, and suppose
$a\le G(x)\le b$ for every $x\in[a,b]$. Let $a<x_L\le x_H<b$ and suppose
\[
  G(x)<x\ \text{ for all }x\in(a,x_L],
  \qquad
  x<G(x)\ \text{ for all }x\in[x_H,b).
\]
No assumption is made about the sign of $G(x)-x$ on $(x_L,x_H)$, and $G$ is not
assumed monotone anywhere. Put
\[
  t=\inf\{\,x\in[x_L,x_H]:x\le G(x)\,\},
  \qquad
  t'=\sup\{\,x\in[x_L,x_H]:G(x)\le x\,\}.
\]
Then $x_L<t\le t'<x_H$, both are fixed points of $G$, and $G(x)<x$ for every
$x\in[x_L,t)$ while $x<G(x)$ for every $x\in(t',x_H]$. Every orbit of $G$ started
in $[a,t)$ converges to $a$; every orbit started in $(t',b]$ converges to $b$;
and $t=t'$ if and only if $G$ has exactly one fixed point in $[x_L,x_H]$.
\end{proposition}

The construction is the whole proof. Continuity makes
$\{x\in[x_L,x_H]:x\le G(x)\}$ closed, so its infimum is attained and lies in it,
giving $t\le G(t)$; and every state below an infimum is outside the set, which
\emph{is} the statement $G(x)<x$ for $x<t$, whence $G(t)\le t$ by continuity from
the left. Because every fixed point of the window lies in the set, $t$ is the
least one. The mirror argument gives $t'$, and $t\le t'$ because $t$ is itself a
fixed point and therefore a member of the second set.

Two features are load-bearing rather than incidental, and both are settled by
counterexample rather than left to the reader's confidence.

\emph{Continuity cannot be weakened.} A map jumping from $0$ to $1$ at the middle
of a window satisfies both endpoint conditions and has no fixed point in the
window at all. Everything the proposition delivers is carried by continuity.

\emph{The interval $[t,t']$ can be non-degenerate.} Take the window $[-1,2]$
itself and the map $x\mapsto\min\{x,2x\}+\max\{0,x-1\}$, which is continuous
there, falls strictly at $-1$, rises strictly at $2$, and fixes every point of
$[0,1]$. Its two derived crossings over that window are $0$ and $1$. So the
inequality $t\le t'$ admits no upgrade to equality, and the region between the
crossings is a genuine gap in the conclusion rather than an artifact of how it is
proved.

The orbit statements do not use monotonicity of $G$. Once $G$ lies strictly below
the diagonal throughout an interval, each iterate steps downward because of the
sign, not because $G$ preserves order; the orbit is then decreasing and bounded,
its limit is a fixed point by continuity, and the strict sign forces that limit to
the endpoint. This is what allows the proposition to be applied across a region
where the sanction rule is known to be continuous but is not known to be monotone,
which is the situation between two certified basins.

\subsection{General closed-loop invariant laws and the affine specialization}

Let $g:[0,1]\to\mathbb R$ be a continuous closed-loop drift and let
$T_{g,h}\mu$ be the law of
\[
x'=\operatorname{clip}\bigl(g(x)+\eta\bigr),
\qquad x\sim\mu,\quad \eta\sim\operatorname{Unif}[-h,h],
\]
with $h>0$ and independent noise.  The formal result is stated at this
general drift level.  Continuity of $g$ is enough for existence of an
invariant probability law.  If, in addition, there is an $L\in(0,1)$
such that
\[
|g(x)-g(y)|\le L|x-y|
\quad\text{for all }x,y\in[0,1],
\]
then
\[
d(T_{g,h}\mu,T_{g,h}\nu)\le L\,d(\mu,\nu),
\qquad
d(T_{g,h}^{\,n}\mu,\lambda)\le L^n d(\mu,\lambda),
\]
and $\lambda$ is the unique invariant probability law.  This
contraction condition is not an additional maintained assumption of the
economic model; it identifies the conditional uniqueness-and-attraction
region of the already general transition framework.

The proof does not use the \citet{HopenhaynPrescott1992} stochastic-order
argument, which needs an order on measures the current libraries lack.
Instead, it constructs the dual Kantorovich--Rubinstein metric
\citep{KantorovichRubinstein1958}, in the duality form of
\citet[\S11.8]{Dudley2002},
$d(\mu,\nu)=\sup_{f\ 1\text{-Lipschitz}}
(\int f\,d\mu-\int f\,d\nu)$, and proves that the closed-loop adjoint maps
$1$-Lipschitz tests to $L$-Lipschitz tests.  Uniqueness is genuine
measure equality, not merely $d=0$: Lipschitz ramps decreasing to
interval indicators show that equal integrals against every Lipschitz
test imply equality of the two probability measures.

The transfer operator and its adjoint identity are constructed directly.
$T_{g,h}\mu$ is the pushforward of the product of $\mu$ and the box-noise
law under the closed-loop-plus-noise map, and
\[
\int f\,d(T_{g,h}\mu)
=\int\!\left[\operatorname{movAvgReal}\,h\,f\bigl(g(x)\bigr)\right]d\mu(x)
\]
follows from the change-of-variables formula and the Bochner Fubini
theorem.  Existence is obtained by the Krylov--Bogolyubov /
Ces\`aro-average argument
(\citealp{KrylovBogolyubov1937}; textbook form
\citealp[Thm.~12.10]{StokeyLucas1989}).  On the compact state space the
Ces\`aro averages have a weak cluster point; the Feller adjoint identity
makes that point invariant against every bounded continuous test, and
the determining class upgrades this to measure-level fixed-point
equality.

The affine drift is a specialization, not the primitive theorem.
For $g_\delta(x)=(1-\delta)x+\delta$, the drift modulus is
$L=1-\delta$, the general operator is definitionally the affine box
transfer, and the preceding bounds become
\[
d(T_{\delta,h}\mu,T_{\delta,h}\nu)
\le(1-\delta)d(\mu,\nu),
\qquad
d(T_{\delta,h}^{\,n}\mu,\lambda)
\le(1-\delta)^n d(\mu,\lambda).
\]
Equivalently, its dual pullback
$(Sf)(x)=\operatorname{movAvgReal}\,h\,f((1-\delta)x+\delta)$ is
$(1-\delta)$-Lipschitz whenever $f$ is $1$-Lipschitz.  Thus the chain with the familiar
affine drift inherits the general theorem's unique, globally
attracting invariant distribution.

The same construction covers the \emph{endogenous} general-$\F$ chain
by taking $g(x)=\F(x,\sigma_{\mathrm{real}}(x))$ whenever this
equilibrium-induced drift is continuous.  Existence then holds with no
slope restriction.  What does not follow automatically is uniqueness:
that conclusion requires the displayed contraction bound.  At the
affine example economies with a positive stage slope the closed-loop factor is exactly
$(1-\delta)+\delta\cdot\tfrac98=1+\delta/8>1$: the slope $\tfrac98$ of the
sanction rule exceeds one precisely because the state--sanction cross-partial
dominates the sender's concavity, the same
$\tfrac{27}{20}$-vs-$\tfrac65$ frontier that closes the
distributional-equilibrium route.  The contraction route is therefore
provably unavailable there, while existence remains.  Economically this
is the appropriate boundary: an equilibrium chain with two invariant
basins need not have a unique globally attracting law. What is
established is only that this sufficient route to such a law is
unavailable. Neither multiplicity of stationary laws nor metastability
is claimed, and a blocked sufficient condition is not evidence for
either.

\subsection{Construction of the shocked (purified) existence (figure)}

\begin{center}
\centering
\resizebox{0.70\textwidth}{!}{%
\begin{tikzpicture}[
  >={Stealth[length=2.2mm]},
  every node/.style={font=\scriptsize},
  box/.style={draw, rounded corners=2pt, align=center, inner sep=3pt,
              text width=33mm, minimum height=8mm},
  prim/.style={box, fill=black!7, draw=black!45},
  der/.style={box, fill=green!12, draw=green!50!black},
  ok/.style={box, fill=teal!14, draw=teal!65},
  thmok/.style={box, fill=orange!22, draw=orange!80!black, very thick},
  todo/.style={box, fill=red!7, draw=red!60, dashed, text width=33mm},
  edge/.style={->, thick, draw=black!72},
  wk/.style={->, thick, dashed, draw=red!65},
]
% main construction chain
\node[prim]  (K)  at (0,9.2)   {\textbf{He--Sun shock kernel} $\kappa$\\ atomless probability measure on the transition noise $\eta$};
\node[der]   (SM) at (0,7.0)   {\textbf{smear} operator\\ $V\mapsto\!\int V(\mathrm{clip}(\cdot+\eta))\,\kappa(\mathrm d\eta)$\\[1pt] {\scriptsize nonexpansive $\checkmark$ (from $\int\kappa=1$)}};
\node[ok]    (BS) at (-3.3,4.7) {Shocked \textbf{Sender} Bellman\\ $\beta_S$-contraction $\checkmark$\\[1pt] $\Rightarrow \Vs^{\mathrm{shock}}$ (Banach)};
\node[ok]    (BT) at (3.3,4.7)  {Shocked \textbf{Target} Bellman\\ $\beta_T$-contraction $\checkmark$\\[1pt] $\Rightarrow \Vt^{\mathrm{dyn,shock}}$, Bellman eq $\checkmark$};
\node[ok]    (EQ) at (0,2.6)    {\textbf{Shocked equilibrium object}\\ {\scriptsize both shocked Bellman fixed-point eqs}};
\node[thmok] (CO) at (-3.3,0.4) {\textbf{Corner} existence $\checkmark$\\ {\scriptsize from the kernel}};
\node[thmok] (IN) at (3.3,0.4)  {\textbf{Interior} existence $\checkmark$\\ {\scriptsize state-varying leader; slaved op contraction $\beta_S+L_{U_S\!,a}L_\alpha<1$}};
% lead_lip reduction column (right)
\node[der]   (SL) at (7.0,7.0)  {\textbf{smear\_lipschitz} $\checkmark$\\ {\scriptsize uniform-kernel moving-average bound}};
\node[der]   (AML) at (7.0,4.4) {\textbf{argmax-Lipschitz core} $\checkmark$\\ {\scriptsize strong concavity $+$ action-Lip.\ objective-diff $\Rightarrow$ Lipschitz argmax}};
\node[der]   (LL) at (7.0,2.2)  {\textbf{lead\_lip} $\checkmark$\\ {\scriptsize argmax-core, \emph{reduced to} the Stackelberg coupling}};
\node[todo]  (SK) at (7.0,-0.4) {\textbf{Reusable-theorem inputs}\\ {\scriptsize general $h$: three inputs; abstract $h=1$: leader strong concavity $+$ coupling; exact $P_3$, $h=1$: bypassed}};
% edges
\draw[edge] (K) -- (SM);
\draw[edge] (SM) -- (BS);  \draw[edge] (SM) -- (BT);
\draw[edge] (BS) -- (EQ);  \draw[edge] (BT) -- (EQ);
\draw[edge] (EQ) -- (CO);  \draw[edge] (EQ) -- (IN);
\draw[edge] (K.west) to[out=180,in=160] node[pos=0.5,left=0pt,font=\tiny,align=center]{nonexp.\\only} (CO.west);
\draw[edge] (SM) -- (SL);
\draw[edge] (SL) -- (AML);
\draw[edge] (AML) -- (LL);
\draw[edge] (LL) -- (IN) node[midway,above,font=\tiny]{discharged};
\draw[wk]   (SK) -- (LL);
\end{tikzpicture}}
\captionof{figure}{\textbf{Construction of the shocked (purified) existence.}
Solid/colored boxes are established results
($\checkmark$); the dashed red box records the reusable general interior
theorem's remaining analytic hypotheses. From the
He--Sun atomless kernel $\kappa$ (grey) the smoothing operator $\mathrm{smear}$
is derived (green), with its nonexpansiveness proved from $\int\kappa=1$. That
nonexpansiveness alone makes the shocked sender and target Bellman operators
$\beta_S$- and $\beta_T$-contractions (teal), whose Banach fixed points are the
shocked value functions and satisfy the two shocked Bellman equations; these
assemble the shocked-equilibrium object. Two existence theorems (orange) follow:
the \emph{corner} equilibrium directly from the kernel, and the \emph{interior}
(state-varying-leader) equilibrium from the slaved-operator contraction under the
condition $\beta_S+L_{U_S\!,a}L_\alpha<1$. The interior case's input
$\mathrm{lead\_lip}$ (leader argmax Lipschitz in $\Vs$) is the argmax-Lipschitz
core applied to the leader objective \emph{difference}; the $\sigma$-Lipschitz
and $\Vt^{\mathrm{dyn}}$-Lipschitz bounds (from $\mathrm{smear\_lipschitz}$,
discharged from the uniform kernel by the moving-average window-shift estimate)
bound the \emph{magnitude} of that difference, but its variation in the leader's
own action is the consumed \emph{Stackelberg-coupling} estimate (part of the
dashed-red package), so $\mathrm{lead\_lip}$ holds given that estimate and
composite strong concavity.  Joint follower continuity and strong concavity also
derive the canonical leader selector, including its feasibility, continuity,
and pointwise optimality; at $h=1$, exact affine smoothing supplies the former
condition for every bounded-continuous continuation. Separately, the exact
$P_3$, $h=1$ construction solves the two-spread Bellman system and verifies the
equilibrium directly, bypassing the dashed-red inputs without claiming them as
general lemmas. The coupling is
genuine: because the sender's response $\sigma(x,a)$ enters the leader objective
and depends on the leader's own action $a$, the action-direction estimate is not
supplied by the shock alone, and a reference sweep confirms it is not derivable
from a smoother kernel (kernel smoothing regularizes the state direction only;
the feedback-Stackelberg literature also assumes the analogous leader-side
regularity). It is isolated as one estimate, exactly as the deterministic
development isolates strong concavity. The deterministic model is the
$\kappa\to\delta_0$ limit, where $\mathrm{smear}=\mathrm{id}$ and the interior
route is blocked.}
\label{fig:shock-construction}
\end{center}
\clearpage
